\section{Introduction}
\label{sec:introduction}

Many failures of language models begin at an identifiable token. A harmful
response has a first harmful instruction, a hallucinated answer has a first
unsupported claim \citep{niu2024ragtruth}, a wrong solution has a first
erroneous step \citep{zheng2025processbench}, and a degenerate response has a
point where it falls into a repetition loop that runs until the token budget is
exhausted \citep{holtzman2020curious,welleck2020neural,pipis2025wait}. A monitor
of such failures can act at three moments. A \emph{response classifier} judges
the finished response, a \emph{failure detector} flags the failure once it is
under way, and an \emph{early-warning monitor} fires before it starts. Only the
last leaves time to stop the generation, steer it
\citep{lee2025cast,hedstrom2025mera} or roll it back and resample
\citep{zhang2025backtracking} before the failure is produced. This paper studies
how to evaluate early-warning monitors, using repetitive degeneration as the
case study (\Cref{fig:degeneration_example}).

Linear probes on hidden states are a cheap way to build such monitors. They
detect deception \citep{goldowskydill2025deception}, high-stakes interactions
\citep{mckenzie2025highstakes}, hallucinated entities
\citep{obeso2025realtime}, cyber-offensive misuse in a production system
\citep{kramar2026gemini}, and repetition loops
\citep{xie2025wordsalad,yu2025breaking,duan2026circular}. Most of this work
evaluates the probe as a classifier. Deception, high-stakes and hallucination
probes report AUROC per response or per entity, together with recall at a low
false-positive rate on benign data
\citep{goldowskydill2025deception,mckenzie2025highstakes,obeso2025realtime}.
Degeneration probes report accuracy and AUROC per chunk or per sentence
\citep{xie2025wordsalad,duan2026circular}. These metrics measure whether a
response or segment is labelled correctly. When every segment after the onset
counts as positive, a probe that fires only once the failure is visible scores
as well as one that warns before it. Some work does measure timing, but not
against an annotated onset under a fixed false-alarm budget. Streaming
harmful-content monitors report the share of the response read before the
monitor stops it \citep{li2025fineharm} or whether the first flag falls inside
the annotated unsafe sentence \citep{qwen2025qwen3guard}.
\citet{duan2026circular} report lead times for loop prediction on balanced
test sets. \citet{itkin2026quickest} argue that token-level AUROC hides how fast
a hallucination detector reacts and measure the delay after the onset at a
fixed false-alarm rate, which gives no credit to an alarm raised before the
onset.

Degeneration makes the problem concrete. We sample responses from
\textsc{Apertus-8B-Instruct} \citep{apertus2025},
\textsc{Llama-3.1-8B-Instruct} \citep{grattafiori2024llama} and
\textsc{Mistral-7B-Instruct} \citep{jiang2023mistral} with one sampling
configuration shared by all three models (temperature $0.7$, top-$p$ $0.9$, up
to $4{,}096$ new tokens). It is not the recommended default of each model
(\todo{list each model's recommended setting in \Cref{app:data}}), but it is a
stochastic setting of the kind used in deployment, unlike the greedy or
near-zero-temperature decoding of prior degeneration work
\citep{xie2025wordsalad,duan2026circular}, and sharing it removes decoding as a
source of difference between models. We do not balance the classes, so
degeneration occurs at its natural rate under this setting, $3.0\%$, $7.5\%$
and $10.7\%$ of validation responses. A monitor that never alarms then reaches $97.0\%$, $92.5\%$ and
$89.3\%$ accuracy. A degenerate response runs to the token limit, so a rule
that reads nothing from the model and flags only responses that reach the limit
reaches an AUROC of \todo{1.000} on our test split. Response-level metrics
cannot tell this rule from a probe that warns hundreds of tokens before the
onset.

We therefore borrow the evaluation used where early warnings are already
deployed. Seizure prediction, quickest change detection and clinical early
warning count alarms relative to the moment the event starts, with the rate of
false alarms held fixed \citep{winterhalder2003seizure,basseville1993detection,hyland2020circulatory}.
\method{} (\emph{onset-aligned warning diagnostics}) applies this to token-level
monitors of language models. It fixes the share $\beta$ of healthy responses
allowed to raise an alarm, and reports, relative to the annotated onset, how much
of the degeneration a monitor flags, how much of the window just before the onset
it flags, how often its first alarm precedes the onset, and by how many tokens.
To our knowledge, no prior evaluation of language-model monitors credits alarms
raised before an annotated onset under a fixed false-alarm budget.

Applied to linear probes on the residual stream of all three models, \method{}
changes which probe one would select. AUROC places every trained probe within
\todo{0.004} of its maximum, and ranking probes by AUROC selects probes with
\todo{3.9}$\times$ lower warning coverage than ranking by \method{}
(\S\ref{sec:rq-metrics}). A similar gap has been reported for clinical
early-warning models, whose discrimination falls to chance once only
predictions made before treatment starts are counted
\citep{kamran2024sepsis}. The probes do carry an early-warning signal. At a
$5\%$ false-alarm budget their first alarm precedes the onset on
\todo{60--73}\% of the degenerate responses they flag, on a test split with
human-verified onsets (\S\ref{sec:rq-signal}). Our contributions are:
\begin{itemize}
  \item \method{}, an evaluation of token-level monitors relative to an
  annotated failure onset under a response-level false-alarm budget, adapted
  from event-based early-warning evaluation (\S\ref{sec:onward}).
  \item Evidence that AUROC and accuracy, computed per response, saturate on
  this task and rank probes against their early-warning behaviour, across
  training strategies and layers (\S\ref{sec:rq-metrics}).
  \item Linear probes for \textsc{Apertus}, \textsc{Llama} and \textsc{Mistral}
  that warn before the onset at a $5\%$ false-alarm budget, and the finding that
  low-rank adaptation of the layers below the probe more than doubles their
  warning coverage (\S\ref{sec:rq-signal}, \S\ref{sec:rq-adaptation}).
  \item A release of these probes with thresholds calibrated at $1\%$, $5\%$
  and $10\%$ false-alarm budgets.
\end{itemize}

\method{} needs only a token-level scorer, a label for where the failure starts,
and a set of healthy responses on which to fix the false-alarm budget. These
already exist for other failures. Harmful content has been annotated at the
token level \citep{li2025fineharm} and with the position where the unsafe part
of a response begins \citep{singteam2026singprobe}, and hallucinations have been
annotated at the word level \citep{niu2024ragtruth}. The right budget depends on
what a false alarm costs. For degeneration it costs stopping or steering a
healthy response, which is why we take $\beta=5\%$ as our working point. Where a
missed failure is the larger cost, a different budget will be preferable.
\Cref{app:related} discusses related work in full.
